\chapter{The core model}
\label{ch:core}

The core model is the simplest model that is consistent with how oceanic residual depths are measured.
It is what the first inversion (Chapter~\ref{ch:single}) uses, and every later stage reduces to it in a stated limit.
If you can derive this chapter on a whiteboard, you can defend the first half of the thesis chapter.

\section{The three assumptions}

\begin{enumerate}[nosep]
\item \textbf{The lithosphere has no strength.}
It rides on the top of the flow and is lifted or lowered as a whole by $w$, with no bending stresses.
\item \textbf{The ocean covers the whole surface and follows the geoid.}
Sea level is an equipotential, so where the geoid rises by $N$ the water deepens by $N$.
\item \textbf{The mass moved by the deflection sits at the surface.}
When a column of lithosphere is lifted by $w$, asthenosphere of density $\rho_{\mathrm{LAB}}$ (PREM's density at the top of the flow, the material immediately below the column) flows into the gap at its base, and the column's top pushes water aside.
The column itself is unchanged, so the mass above any fixed depth changes by $+\rho_{\mathrm{LAB}}w$ at the base and $-\rho_w w$ at the top.
For any column, lifted rigidly, the changes at all its density contrasts telescope: the surface gives $(\rho_{\mathrm{top}}-\rho_w)w$, each internal contrast $(\rho_{\mathrm{below}}-\rho_{\mathrm{above}})w$, the base $(\rho_{\mathrm{LAB}}-\rho_{\mathrm{bottom}})w$, and the sum is always $(\rho_{\mathrm{LAB}}-\rho_w)w$.
So the restoring force does not depend on the column's internal structure; only where the mass sits does, and that only matters for the geoid.
``Density $\rho_{\mathrm{LAB}}$ throughout'' is the model's shorthand for where the moved mass goes, not a claim about real lithosphere: residual depths have had the cooling plate, sediments and crust removed, so the data treat the column's structure as accounted for.
For real oceanic lithosphere it is close: the large contrasts are at the seafloor and at the Moho, within about \qty{7}{km} of the surface, while the lid grades into the asthenosphere.
If the column's density is $\rho_{\mathrm{LAB}}$ throughout, the base swap is lithosphere for asthenosphere of the same density and changes nothing, so the whole change, $(\rho_{\mathrm{LAB}}-\rho_w)\,w$ per unit area, sits at the top of the column.
We put it as a thin sheet at $r=R$, the undeflected top of the column; the mass really occupies $R$ to $R+w$, but the difference is second order.
Keeping a separate sheet at the base would not help: for the uniform column it is zero. The tempting split, $+\rho_{\mathrm{LAB}}w$ at the base and $-\rho_w w$ at the top, is wrong: it forgets that the column itself moved, lifting in effect an empty column, and gives a spurious mass dipole in the geoid.
(The real seafloor is \qtyrange{2}{6}{km} below sea level; putting the sheet at $R$ rather than there changes its geoid by under \qty{1}{\percent} at $\ell\le10$.)
\end{enumerate}

\begin{decision}{Why these three, for oceanic residual depths}
Each assumption matches what the residual-depth data already assume.
\begin{itemize}[nosep]
\item Residual depths are \emph{water-loaded}: the depth anomaly is what it would be if the deflection were filled with water.
That is assumption 2 at sea.
\item Residual depths remove the plate's cooling, sediments and crust, i.e.\ the lithosphere's own structure, leaving a uniform column that only moves up and down: assumption 1 for the load, and assumption 3 for where its mass goes.
\item The flexural strength of the oceanic plate matters only at wavelengths shorter than those we use: for $T_e=\qty{40}{km}$ flexure changes the deflection by about \qty{1}{\percent} at degree 30, \qty{3.5}{\percent} at degree 40, and negligibly below degree 20 (Chapter~\ref{ch:plate}).
\end{itemize}
So the core model is not a crude version of the full model: for this observable it \emph{is} the model, and the extra parts of the full model change the prediction by \qtyrange{1}{4}{\percent} (below).
\end{decision}

\section{The force balance}

Consider a column of unit area at the top of the flow.
The mantle pushes up on it with the load $q$ (from Chapter~\ref{ch:flow}).
Lifting the column by $w$ costs weight in two ways.
First, by Archimedes: the column now sits $w$ less deep in the mantle, losing buoyancy $\rho_{\mathrm{LAB}}\,g\,w$, and pushes $w$ further into the water, gaining buoyancy $\rho_w g\,w$.
Equivalently, the mass above a fixed depth below the column has increased by $(\rho_{\mathrm{LAB}}-\rho_w)\,w$: dense rock is drawn in below and only light water is pushed aside above.
The restoring force is therefore $(\rho_{\mathrm{LAB}}-\rho_w)\,g\,w$.
Second, the sea surface has risen by $N$ (it follows the geoid), so the column carries an extra layer of water of thickness $N$, with weight $\rho_w g_0 N$.
Force up = force down:
\begin{equation}
    q = (\rho_{\mathrm{LAB}}-\rho_w)\,g\,w + \rho_w\,g_0\,N .
    \label{eq:c-force}
\end{equation}
Here $g$ is flowkernels' gravity at the top of the flow (\qty{9.91}{m/s^2}; integrating PREM gives about \qty{9.88}{m/s^2}) and $g_0$ at the surface (\qty{9.82}{m/s^2}).
Strictly the water gained at the top sits at the surface, so its share of the restoring force should use $g_0$; we follow the plate code and use one $g$ for the column, an error of about \qty{0.4}{\percent} in the restoring force.
The full model also weights the geoid-following water with $g$ rather than $g_0$, so the core model matches the full model's no-plate, all-ocean limit only to about $10^{-3}$, not exactly.
None of these differences matters at the level of the data errors, but do not claim they are exact.

\section{The geoid: four sources}

The geoid at the surface is the sum of the potentials of every mass the model moves, divided by $-g_0$.
There are four.
\begin{enumerate}[nosep]
\item \textbf{The density anomaly} in the mantle, $N^{\rho}$.
\item \textbf{The deflected core--mantle boundary}, $N^{\mathrm{cmb}}$: the anomaly's flow pushes the CMB up or down, and the deflection is a sheet of mass $(\rho_{\mathrm{core}}-\rho_{\mathrm{mantle}})\,u_r$.
\item \textbf{The moved lithospheric column}, a sheet of mass $(\rho_{\mathrm{LAB}}-\rho_w)\,w$ at the surface.
\item \textbf{The water that follows the geoid}, a sheet $\rho_w N$ at the surface.
\end{enumerate}
The first two together are $D = N^{\rho}+N^{\mathrm{cmb}}$, from the kernel $K_N$ of Chapter~\ref{ch:flow}.

\paragraph{The geoid of a surface sheet.}
A thin sheet of surface density $\sigma_{\ell m}$ at radius $r$ has, outside it, the potential
\begin{equation}
    \Phi_{\ell m}(R) = -\frac{4\pi G r}{2\ell+1}\left(\frac{r}{R}\right)^{\ell+1}\sigma_{\ell m},
\end{equation}
which follows from solving Poisson's equation with a delta-function source and matching the potential and the jump in its gradient across the sheet.
Dividing by $-g_0$ gives the geoid per unit surface density:
\begin{equation}
    S_\ell(r) = \frac{4\pi G\,r}{(2\ell+1)\,g_0}\left(\frac{r}{R}\right)^{\ell+1}, \qquad
    S_\ell(R) = \frac{4\pi G R}{(2\ell+1)\,g_0}.
\end{equation}
At $\ell=2$, $S_2(R)\approx \qty{1.1e-4}{m}$ per \si{kg/m^2}: a degree-2 layer of water \qty{1}{m} thick ($\qty{1000}{kg/m^2}$) raises the geoid by about \qty{0.11}{m}.
The water sheet itself is therefore not small; but because the column sinks under the water's weight, its net effect on $N$ nearly cancels (next section).

The total geoid is
\begin{equation}
    N = D + S_\ell(R)\left[(\rho_{\mathrm{LAB}}-\rho_w)\,w + \rho_w N\right] .
    \label{eq:c-geoid}
\end{equation}

\section{The closed-form solution}

Equations~\eqref{eq:c-force} and~\eqref{eq:c-geoid} are two linear equations in two unknowns, $w$ and $N$, for each $(\ell,m)$.
From Equation~\eqref{eq:c-force}, $(\rho_{\mathrm{LAB}}-\rho_w)\,w = (q-\rho_w g_0 N)/g$.
Substituting into Equation~\eqref{eq:c-geoid}:
\begin{align}
    N &= D + S\left[\frac{q-\rho_w g_0 N}{g} + \rho_w N\right]
       = D + \frac{S q}{g} + \rho_w S N\left(1-\frac{g_0}{g}\right), \\
    N &= \frac{D + S q/g}{1 + \rho_w S\,(g_0/g - 1)} ,
    \qquad
    w = \frac{q - \rho_w g_0 N}{(\rho_{\mathrm{LAB}}-\rho_w)\,g} .
\end{align}
The denominator differs from 1 by $\rho_w S (g_0/g-1) \approx 10^{-3}$ at degree 2: the water's extra weight and its extra attraction almost exactly cancel, because when $g=g_0$ the column sinks under the extra water by just enough to keep the total mass at the surface unchanged.
The water still matters, through the term $\rho_w g_0 N$ in the force balance, which changes $w$, and through $h=w-N$.

\paragraph{What is observed.}
Topography is measured from sea level, which is the geoid.
The seafloor moves up by $w$ and sea level by $N$, so the observed depth anomaly is
\begin{equation}
    h = w - N .
\end{equation}
This matters most at low degree, where $N$ is a few per cent of $w$.

\paragraph{Degrees 0 and 1.}
Degree 0 is excluded: a uniform change of density changes the Earth's mass and radius, not its shape, and the residual depths are defined up to a constant, which the inversion fits separately ($b_0$).
At degree 1 we are in the centre-of-mass frame (Chapter~\ref{ch:flow}), so $N_1 = 0$ and
\begin{equation}
    w_1 = \frac{q_1}{(\rho_{\mathrm{LAB}}-\rho_w)\,g}, \qquad h_1 = w_1 .
\end{equation}

\section{Evidence that the core model is right}

\begin{defence}
\Q{Why put the moved mass at the surface rather than at \qty{105}{km}?}
\A Because that is where it is.
Lifting a column by $w$ swaps rock for water at the top (the surface) and swaps lithosphere for asthenosphere at the base; if the column's density equals $\rho_{\mathrm{LAB}}$ throughout (which is what removing the lithosphere's own structure from the data amounts to), the base swap is zero and all of the change is at the top.
Putting it at \qty{105}{km}, as flowkernels' own geoid does, changed the low-degree geoid by factors of 0.2 to 3, because the compensated geoid is a small difference of large terms.

\Q{How do you know the simplifications don't matter?}
\A We compared the core model with the full model (elastic plate, real ocean, lithosphere about WINTERC-G) at the oceanic data points for the same mantle.
The predicted topographies correlate at $\ge 0.996$ and differ in amplitude by \qtyrange{1}{4}{\percent}; the geoids agree within about \qty{15}{\percent} at degrees 2--8, except near the sign change of the geoid kernel where both are small.
The data errors are much larger than these differences.

\Q{Why not use continental data with the core model?}
\A Because on land there is no water in the deflection: the restoring force is $\rho_{\mathrm{LAB}}\,g\,w$, not $(\rho_{\mathrm{LAB}}-\rho_w)\,g\,w$.
The water-loaded prediction overestimates the land deflection by $\rho_{\mathrm{LAB}}/(\rho_{\mathrm{LAB}}-\rho_w)\approx 1.44$.
That is the problem the elastic plate on land and sea solves (Chapter~\ref{ch:plate}).

\Q{Is the ocean term worth keeping?}
\A Yes, though not for the reason one might first think.
The water sheet is about \qty{7}{\percent} of the mantle geoid at degree 2, but its weight sinks the column by just enough to cancel almost all of that: its net effect on $N$ is about \qty{0.1}{\percent}.
Its real effect is on the deflection: the extra water changes $w$ by $-\rho_w N/(\rho_{\mathrm{LAB}}-\rho_w)\approx -0.44N$, about \qty{2}{\percent} of $w$ at degree 2, and sea level enters the observable through $h=w-N$.
\end{defence}

\begin{incode}{\texttt{scripts/single\_scaling/build\_core.py}}
\begin{enumerate}[nosep]
\item Builds \texttt{MantleFlow} with the chosen viscosity, \texttt{lab\_depth=105e3}, \texttt{lmax=40}, and reads its kernels.
\item Reads the tomography (dlnVs in per cent) and evaluates it on the kernel grid (\texttt{tomography\_at}, degree 0 zeroed, divided by 100), and the reference density $\rho_0$ (\texttt{rho0\_on\_kernel\_grid}).
\item For each of the 28 depth nodes of $R(z)$, forms the density $\delta\rho = \rho_0\,B_i(z)\,\mathrm{dlnV_s}$ with the hat function $B_i$, and the load and direct geoid
\begin{lstlisting}
q[:, l, :l+1] = tensordot(K_q[l-1], f[:, :, l, :l+1], axes=(0, 0))
D[:, l, :l+1] = tensordot(K_N[l-1], f[:, :, l, :l+1], axes=(0, 0))
\end{lstlisting}
where \texttt{f} already contains the quadrature weights.
\item Applies the closed form:
\begin{lstlisting}
S = 4*pi*G*R / ((2*l + 1)*g0)
N = (D + S*q/g_top) / (1 + rho_w*S*(g0/g_top - 1))
w = (q - rho_w*g0*N) / ((rho_lab - rho_w)*g_top)
N[:, 1, :] = 0;  w[:, 1, :] = q[:, 1, :] / ((rho_lab - rho_w)*g_top)
h = w - N
\end{lstlisting}
\item Evaluates $h$ at every data point with orthonormal real harmonics (\texttt{harmonics\_at}, Condon--Shortley phase $+1$), giving a column of \texttt{G\_points}.
\item Converts $N$ to the convention of the geoid product: multiplies order $m$ by $(-1)^m$ (Condon--Shortley phase $-1$) and by $(R/a)^{\ell-1}$ to refer it to the product's radius $a$, giving a column of \texttt{G\_geoid}.
\end{enumerate}
Each column is the prediction for $R(z) = B_i(z)$; the full prediction is \texttt{G\_points @ R}.
\end{incode}
